\scanpage{230}
par $Z$ de la partie de $R$, stable sous $Z$,
\deleted{\unsure{formée de|définie}} \inserted{définie} par
\diagram{\begin{tikzcd}[ampersand replacement=\&,row sep=huge]
{\text{(39)}} \& {\deleted{\unsure{Z\in|Z}}\ R(s^{\unsure{\leftrightharpoons|\natural|\leftrightarrows|\flat}},s^{\vee\leftrightharpoons})=
\left\{(s,s^\vee)\in R\ \middle|\ \begin{aligned}&s\text{ au dessus de }s^\leftrightharpoons\\
&s^\vee\text{ au dessus de }s^{\vee\leftrightharpoons}\end{aligned}\right\}}\\
\& { \unsure{E|\mathcal E}(s_0,s_0^\vee)\overset{\mathrm{déf}}=
\left\{(\lambda,\mu)\in Z\times Z\ \middle|\
\begin{aligned}&(\lambda s_0,\mu s_0^\vee)\in R\quad\text{i.e.}\\
&((\mu^{-1}\lambda)s_0,s_0^\vee)\in R\quad\text{i.e.}\\
&\mu^{-1}\lambda s_0\in\underbrace{E(s_0^\vee)}_{\substack{\text{étoile dans }S\\
\text{définie par }s_0^\vee}}\end{aligned}\right\}.}
\arrow[u,"\wr" description,"{(\lambda,\mu)\longmapsto(\lambda s_0,\mu s_0^\vee)}"']
\end{tikzcd}}
On trouve ainsi une bijection can.
\[
\begin{aligned}
 &\text{fibre de }R^\leftrightharpoons\text{ en }(s^\leftrightharpoons,s^{\vee\leftrightharpoons})\simeq
 \deleted{\text{sous-groupe de}}\deleted{\text{fini}}
 \inserted{Z(s_0,s_0^\vee)}\ \inserted{\text{ens. des }}z\in Z\\
 &\hspace{13em}\text{tels que }\ (zs_0,s_0^\vee)\in R\\
 &\hspace{13em}\text{i.e. }\ (s_0,z^{-1}s_0^\vee)\in R\\
 &\hspace{13em}\text{i.e. }\ zs_0\in E
\end{aligned}\tag*{(40)}
\]
Cette bijection en fait, \inserted{qui dépend \inserted{du choix} de $s_0,s_0^\vee$,} ne dépend pas de l'hypothèse
$(s_0\,s_0^\vee)\in R$ -- les deux mbs de (40) \unsure{pouvant|peuvent} être vides,
\deleted{ce qui} L'hypothèse $(s_0,s_0^\vee)\in R$ signifie simplement que
$Z_{s_0,s_0^\vee}\ni1$. Quand on change de $(s_0,s_0^\vee)$, le
sous-ens. en question est remplacé par un translaté.

\deleted{Dans le cas où $Z$-jauge stricte $\varphi_0$, on a l'équivalence}
\inserted{On a l'équivalence [si $Z$ opère librement sur $S,S^\vee$]}
\[
\begin{aligned}
&\deleted{\text{équivalence}}\quad R^\leftrightharpoons\longrightarrow Z^\leftrightharpoons\times Z'{}^\leftrightharpoons
 \text{ injective}\ \Longleftrightarrow\
 \deleted{\unsure{E}}\inserted{\forall} E\in\mathcal E,\ s_0\in E,\\
&\hspace{12em}\lambda\in Z,\ \inserted{\text{on a}}\ \lambdas_0\in E\Longrightarrow\\
&\hspace{12em}\deleted{\lambda s_0\notin E}\quad\lambda=1\\
&\underset{}{\overset{\substack{\text{en termes d'une}\\Z\text{-jauge \uline{stricte}}}}{\Longleftrightarrow}}
 \ \forall(s,t)\in S\times S,\ \inserted{\text{on a }}\{s,t\}\in A_0\
 \deleted{\Longrightarrow}\ ,\ \{s,\lambda t\}\in A_0\Rightarrow\lambda=1\\
&\Longleftrightarrow\ \forall\lambda\in Z,\ \lambda\ne1,\ \text{on a }A_0\cap A_\lambda=\{\unsure{\emptyset}\}\\
&\deleted{\Longleftrightarrow\ \forall\lambda\in Z,\ \unsure{\lambda\ne1},\ A_0\cap A_\lambda=\emptyset
 \quad\text{si on donne une }Z\text{-jauge }\varphi_0,\ \lambda\in Z,\ \lambda\ne1}\\
&\deleted{\Longleftrightarrow\ \{\forall\lambda,\lambda'\in Z,\ \lambda\ne\lambda'\quad
 \text{sont disjointes}\quad A_\lambda\cap A_{\lambda'}=\emptyset}
\end{aligned}\tag*{(41)}
\]
Cette condition signifie que
si $D$ et $D^\vee$ sont \deleted{deux} \inserted{un} rayon et corayon incidents,
alors pour tout $s\in D$, $\exists!s^\vee\in D^\vee$ tel que $s,s^\vee\in D^\vee$.
Si $s$ remplacé par $\lambda s$, $s^\vee$ remplacé par $\lambda s^\vee$. Donc
\deleted{Si (41) \unsure{est tel}} \inserted{(41) équivaut à}
\[
\begin{aligned}
&\text{Pour deux rayons }\deleted{\unsure{D,D^\vee}}\inserted{D\in S^\leftrightharpoons,D^\vee\in S^{\vee\leftrightharpoons}}\text{ incidents,}\\
&\text{la relation }\deleted{\text{entre}}\ D\text{ et }D^\vee\qquad (s,s^\vee)\in R\text{ est}\\
&\text{une relation bijective, établissant \unsure{une isom|un isom|un isom.} de}\\
&Z\text{-torseurs}\\
&\hspace{6em}D\simeq D^\vee
\end{aligned}\tag*{$(41')$}
\]
% archive sheet number: 229

